\documentclass[11pt]{article}
\usepackage[utf8]{inputenc}
\usepackage{amsmath,amssymb}
\usepackage{graphicx}
\usepackage{hyperref}
\usepackage{booktabs}
\usepackage{listings}

\title{On the Convergence Properties of Gradient Descent in High-Dimensional Spaces}
\author{Alice Researcher \and Bob Scientist}
\date{March 2024}

\begin{document}
\maketitle

\begin{abstract}
We study the convergence behavior of gradient descent algorithms applied to
non-convex optimization problems in high-dimensional parameter spaces. Our
analysis reveals that under mild regularity conditions, gradient descent with
appropriate step-size scheduling achieves $\mathcal{O}(1/\sqrt{T})$
convergence to first-order stationary points. We validate our theoretical
findings with experiments on synthetic and real-world datasets.
\end{abstract}

\section{Introduction}\label{sec:intro}

Optimization lies at the heart of modern machine learning \cite{bottou2018}.
Given a loss function $\mathcal{L}: \mathbb{R}^d \to \mathbb{R}$, the goal
is to find parameters $\theta^*$ such that:
\begin{equation}\label{eq:objective}
\theta^* = \arg\min_{\theta \in \mathbb{R}^d} \mathcal{L}(\theta).
\end{equation}

While convex optimization is well-understood \cite{boyd2004,nesterov2004},
the non-convex landscape typical of deep learning presents significant
challenges. Recent work by \citet{allen2019} and \citet{du2019} has made
progress on understanding convergence in over-parameterized regimes.

\paragraph{Contributions} Our main contributions are:
\begin{enumerate}
\item A unified convergence framework for gradient descent variants under
  relaxed smoothness assumptions (Section~\ref{sec:theory}).
\item Sharp convergence rates that improve upon existing bounds by a factor
  of $\log(d)$ (Theorem~\ref{thm:main}).
\item Extensive experimental validation on both synthetic benchmarks and
  standard datasets (Section~\ref{sec:experiments}).
\end{enumerate}

\paragraph{Notation} We write $\|\cdot\|$ for the Euclidean norm,
$\langle \cdot, \cdot \rangle$ for the inner product, and $\nabla f$ for
the gradient of $f$. Constants $c, C > 0$ are universal unless stated
otherwise.

\subsection{Related Work}\label{sec:related}

The study of gradient descent dates back to \citet{cauchy1847}. Modern
analyses can be grouped into several categories:

\begin{description}
\item[Convex optimization] Classical results establish $\mathcal{O}(1/T)$
  convergence for smooth convex functions and $\mathcal{O}(1/\sqrt{T})$ for
  non-smooth variants \cite{nesterov2004}.
\item[Non-convex optimization] Under $L$-smoothness, gradient descent
  achieves $\mathcal{O}(1/\sqrt{T})$ convergence to stationary points
  \cite{nesterov2004}.
\item[Stochastic methods] SGD and its variants introduce noise but maintain
  similar asymptotic rates with appropriate step-size decay
  \cite{robbins1951}.
\end{description}

\section{Theoretical Framework}\label{sec:theory}

\subsection{Assumptions}

We require the following regularity conditions on the objective:

\begin{enumerate}
\item \textbf{Smoothness}: $\mathcal{L}$ is $L$-smooth, i.e.,
  $\|\nabla \mathcal{L}(x) - \nabla \mathcal{L}(y)\| \leq L\|x - y\|$
  for all $x, y \in \mathbb{R}^d$.
\item \textbf{Bounded below}: There exists $\mathcal{L}^* > -\infty$ such
  that $\mathcal{L}(\theta) \geq \mathcal{L}^*$ for all $\theta$.
\item \textbf{Gradient access}: We have access to exact gradients
  $\nabla \mathcal{L}(\theta)$ at any point $\theta$.
\end{enumerate}

\subsection{Main Result}

\begin{equation}
\theta_{t+1} = \theta_t - \eta_t \nabla \mathcal{L}(\theta_t),
\end{equation}
where $\eta_t = \eta / \sqrt{t+1}$ is the step size at iteration $t$.

Our main theorem establishes the convergence rate:

\begin{align}\label{eq:convergence}
\min_{0 \leq t \leq T-1} \|\nabla \mathcal{L}(\theta_t)\|^2
  &\leq \frac{2[\mathcal{L}(\theta_0) - \mathcal{L}^*]}{\eta \sqrt{T}}
     + \eta L \sqrt{T} \\
  &= \mathcal{O}\left(\frac{1}{\sqrt{T}}\right). \nonumber
\end{align}

For the optimal step size $\eta^* = \sqrt{2[\mathcal{L}(\theta_0) -
\mathcal{L}^*] / (LT)}$, we obtain:

\begin{equation*}
\min_{0 \leq t \leq T-1} \|\nabla \mathcal{L}(\theta_t)\|^2
  \leq 2\sqrt{\frac{2L[\mathcal{L}(\theta_0) - \mathcal{L}^*]}{T}}.
\end{equation*}

\subsection{Comparison with Prior Bounds}

Table~\ref{tab:comparison} compares our bounds with prior work.

\begin{tabular}{lccc}
\toprule
Method & Rate & Smoothness & Reference \\
\midrule
GD (convex) & $\mathcal{O}(1/T)$ & $L$-smooth & \cite{nesterov2004} \\
GD (non-convex) & $\mathcal{O}(1/\sqrt{T})$ & $L$-smooth & \cite{nesterov2004} \\
SGD & $\mathcal{O}(1/\sqrt{T})$ & $L$-smooth & \cite{robbins1951} \\
\textbf{Ours} & $\mathcal{O}(1/\sqrt{T})$ & $(L,\alpha)$-smooth & This paper \\
\bottomrule
\end{tabular}

\section{Experiments}\label{sec:experiments}

\subsection{Synthetic Benchmarks}

We evaluate on the \emph{Rosenbrock function}\footnote{Also known as
Rosenbrock's valley or banana function, introduced in 1960.}:
\begin{equation}
f(x, y) = (a - x)^2 + b(y - x^2)^2,
\end{equation}
with standard parameters $a = 1$, $b = 100$.

Results are shown in Figure~\ref{fig:convergence}.

\begin{figure}[htbp]
\centering
\includegraphics{convergence_plot.png}
\caption{Convergence of gradient descent on the Rosenbrock function with
  different step-size schedules. The $x$-axis shows iterations and the
  $y$-axis shows $\log_{10}\|\nabla f\|$.}
\label{fig:convergence}
\end{figure}

\subsection{Implementation Details}

All experiments use the following hyperparameters:

\begin{itemize}
\item Learning rate: $\eta_0 = 0.01$ with cosine annealing
\item Batch size: 256
\item Weight decay: $\lambda = 10^{-4}$
\item Momentum: $\beta = 0.9$
  \begin{itemize}
  \item Nesterov variant used for all non-Adam experiments
  \item Warmup period: 5 epochs
  \end{itemize}
\item Gradient clipping: $\|\nabla\| \leq 1.0$
\end{itemize}

The training loop is implemented as follows:

\begin{verbatim}
for epoch in range(num_epochs):
    for batch in dataloader:
        loss = model(batch)
        loss.backward()
        clip_grad_norm_(model.parameters(), 1.0)
        optimizer.step()
        optimizer.zero_grad()
\end{verbatim}

With language-specific highlighting:

\begin{lstlisting}[language=Python]
def gradient_descent(f, grad_f, x0, lr=0.01, steps=1000):
    """Run gradient descent on function f."""
    x = x0.copy()
    history = [x.copy()]
    for t in range(steps):
        g = grad_f(x)
        x -= lr / np.sqrt(t + 1) * g
        history.append(x.copy())
    return x, history
\end{lstlisting}

\subsection{Results on Standard Datasets}

Table~\ref{tab:results} summarizes our experimental results.

\begin{tabular}{|l|c|c|r|}
\hline
Dataset & Method & Accuracy & Time (s) \\
\hline
MNIST & GD & 98.2\% & 45 \\
MNIST & SGD & 99.1\% & 12 \\
MNIST & Adam & 99.3\% & 15 \\
\hline
CIFAR-10 & GD & 82.1\% & 340 \\
CIFAR-10 & SGD & 93.4\% & 89 \\
CIFAR-10 & Adam & 94.1\% & 102 \\
\hline
\end{tabular}

\section{Discussion}

Our results demonstrate that gradient descent with appropriate step-size
scheduling can achieve near-optimal convergence rates even in
\underline{non-convex settings}. The key insight is that the
\textsc{smoothness} condition can be significantly relaxed.

For further details, visit the project page at
\url{https://example.com/gradient-convergence} or see the companion
repository at \href{https://github.com/example/gradient-descent}{the
GitHub page}.

\begin{quote}
The art of optimization is not in finding the global minimum, but in
knowing when a local minimum is good enough. --- Anonymous
\end{quote}

\begin{quotation}
Gradient descent is remarkably effective in practice, even when theoretical
guarantees are weak. This gap between theory and practice remains one of the
most important open questions in optimization theory. Understanding why
gradient descent works so well in high-dimensional, non-convex landscapes
is crucial for the continued development of machine learning.
\end{quotation}

\section{Conclusion}\label{sec:conclusion}

We have presented improved convergence bounds for gradient descent under
relaxed smoothness assumptions. Our framework unifies several existing
results (see Section~\ref{sec:related}) and provides tighter bounds in
the non-convex case (Equation~\eqref{eq:convergence}).

Future work includes extending these results to the stochastic
setting\footnote{Preliminary results suggest similar rates hold for SGD
with variance reduction.} and investigating connections to neural tangent
kernel theory\footnote{See \cite{jacot2018} for background on NTK.}.

\end{document}
